\documentclass[12pt]{article}
\usepackage[utf8]{inputenc}
\usepackage[LGR,T1]{fontenc}
\usepackage[greek,english]{babel}
\usepackage{alphabeta}
\usepackage{listings}
\usepackage{xcolor}
\usepackage{booktabs}
\usepackage{multirow}
\usepackage{graphicx}
\usepackage{amsmath}
\usepackage{float}
\usepackage[unicode]{hyperref}
\usepackage{subcaption}
\usepackage[a4paper, total={17cm, 25cm}]{geometry}
\usepackage{setspace}
\setstretch{1.15}
\usepackage{tikz}
\usepackage{tikz}
\usepackage{caption}
\usepackage{subcaption}
\usepackage{booktabs}

\begin{document}

\begin{titlepage}
    \begin{center}
        \vspace*{1cm}
        
        % Λογότυπο (Αν έχεις το αρχείο ntua_logo.png, αλλιώς σχολίασε τη γραμμή)
        % \includegraphics[width=0.2\textwidth]{ntua_logo.png}\\[1cm]
        
        \Large
        \textbf{ΕΘΝΙΚΟ ΜΕΤΣΟΒΙΟ ΠΟΛΥΤΕΧΝΕΙΟ}\\
        \large
        ΣΧΟΛΗ ΗΛΕΚΤΡΟΛΟΓΩΝ ΜΗΧΑΝΙΚΩΝ ΚΑΙ ΜΗΧΑΝΙΚΩΝ ΥΠΟΛΟΓΙΣΤΩΝ\\[1.5cm]
        
        \vfill
        
        \Large
        \textbf{Μάθημα: Όραση Υπολογιστών}\\[0.5cm]
        
        \vfill
        
        \huge
        \textbf{Αναφορά: Εργαστήριο 2}\\[2cm]
        
        \vfill
        
        \Large
        \textbf{Στοιχεία Φοιτητών:}\\[0.5cm]
        \large
        Παναγόπουλος Μιχαήλ, Α.Μ: 03122061\\
        Φειδάκης Αθανάσιος, Α.Μ: 03122676\\[2cm]
        
        \vfill
        
        \large
        Αθήνα, Μάρτιος 2026
        
    \end{center}
\end{titlepage}

\newpage

\section*{Μέρος 1: Παρακολούθηση Προσώπου και Χεριών με Χρήση
της Μεθόδου Οπτικής Ροής των Lucas-Kanade}


\subsubsection*{1.1.1 Υλοποίηση του Αλγόριθμου των Lucas-Kanade}

\hspace{0.6cm}Η συνάρτηση \texttt{lk(I1, I2, features, $\rho$, $\varepsilon$, $d_{x0}=None$, $d_{y0}=None$)} υπολογίζει την οπτική ροή από την $I_{n-1}$ προς την $Ι_{n}$. Δεδομένων σημείων ενδιαφέροντος $\mathbf{x}_i$ εντοπισμένων στο τρέχον πλαίσιο $I_n$ (\texttt{I2}), εκτιμά τη μετατόπιση $\mathbf{d}_i$ τέτοια ώστε $I_{n-1}(\mathbf{x}_i + \mathbf{d}_i) \approx I_n(\mathbf{x}_i)$. Με άλλα λόγια προσπαθεί να εντοπίσει την προέλευση των εντοπισμένων στην $I_{n}$ σημείων ενδιαφέροντος από την $Ι_{n-1}$. 

Τα σημεία εντοπίζονται με \texttt{cv2.goodFeaturesToTrack} πάνω στο $I_n$, ο οποίος υλοποιεί τον ανιχνευτή Shi-Tomasi. Χρησιμοποιούνται οι παράμετροι \texttt{maxCorners=30}, \texttt{qualityLevel=0.1} και \texttt{minDistance=5} pixels. Το \texttt{qualityLevel=0.1} επιλέγει σημεία με απόκριση τουλάχιστον 10\% της μέγιστης, φιλτράροντας ασθενείς γωνίες χωρίς να γίνεται υπερβολικά επιλεκτικό. Το \texttt{minDistance=5} εξασφαλίζει χωρική κατανομή των σημείων ώστε να αποφευχθεί η συσσώρευσή τους σε μία μόνο ακμή, δίνοντας πιο αντιπροσωπευτική εκτίμηση μετατόπισης του bounding box. Το \texttt{maxCorners=30} ανά crop περιορίζει τον υπολογιστικό φόρτο, αφού ο LK επιλύει ένα 2×2 σύστημα για κάθε σημείο σε κάθε επανάληψη.

\subsection*{Παρατηρήσεις υλοποίησης - Σχολιασμός Παραμέτρων:}
\begin{enumerate}
    \item \textbf{Παράθυρο Gaussian:} Για κάθε σημείο $\mathbf{x}_i$ ορίζεται ένα τετράγωνο παράθυρο μεγέθους $(2\lceil 3\rho \rceil + 1)^2$ pixels, σταθμισμένο με 2D Gaussian $G_\rho$ που κατασκευάζεται μέσω της συνάρτησης \texttt{cv2.getGaussianKernel}. Η παράμετρος $\rho$ ελέγχει την ακτίνα ολοκλήρωσης του τανυστή δομής. Μικρά $\rho$ εστιάζουν σε τοπικές δομές (κατάλληλο για μικρή κίνηση), μεγάλα $\rho$ ενσωματώνουν ευρύτερη γειτονιά (αντοχή σε θόρυβο, αλλά απώλεια τοπικότητας).

    \item \textbf{Χωρικές Παράγωγοι:} Οι παράγωγοι $I_{1x}, I_{1y}$ υπολογίζονται μία φορά έξω από τον βρόχο με \texttt{cv2.Sobel} στο $I_{n-1}$. Σε κάθε επανάληψη δειγματοληπτούνται στις μετακινούμενες θέσεις $\mathbf{x}_i + \mathbf{d}_i$ μέσω \texttt{map\_coordinates}, υλοποιώντας δηλαδή interpolation για την εύρεση τιμών σε subpixel συντεταγμένες.

    \item \textbf{Συναρμολόγηση Τανυστή Δομής:} Ο πίνακας $M$ και το διάνυσμα $\mathbf{b}$ υπολογίζονται per-feature μέσω σταθμισμένων αθροισμάτων επί του παραθύρου $W(\mathbf{x}_i)$:
    \[
    M = \sum_{\mathbf{x} \in W} G_\rho \begin{pmatrix} I_{1x}^2 & I_{1x}I_{1y} \\ I_{1x}I_{1y} & I_{1y}^2 \end{pmatrix} + \varepsilon I, \quad
    \mathbf{b} = \sum_{\mathbf{x} \in W} G_\rho \begin{pmatrix} I_{1x} \cdot E \\ I_{1y} \cdot E \end{pmatrix}
    \]
    όπου $E = I_n(\mathbf{x}_i) - I_{n-1}(\mathbf{x}_i + \mathbf{d}_i)$ το σφάλμα φωτεινότητας. Η παράμετρος $\varepsilon$ προστίθεται στα διαγώνια στοιχεία ως regularization, αποτρέποντας ιδιόμορφο $M$ σε ομοιογενείς περιοχές. Μικρό $\varepsilon$ (π.χ.\ $10^{-4}$--$10^{-2}$) αφήνει τον $M$ να καθορίζει ελεύθερα τη λύση, ενώ μεγάλο $\varepsilon$ κυριαρχεί στη διαγώνιο και συρρικνώνει τις εκτιμήσεις $(u_x, u_y)$ προς το μηδέν, οδηγώντας σε συστηματική υποεκτίμηση της πραγματικής μετατόπισης.

    \item \textbf{Επίλυση με Κανόνα Cramer:} Το σύστημα $M\mathbf{u} = \mathbf{b}$ επιλύεται αναλυτικά με την μέθοδο του Cramer. Η διόρθωση $\mathbf{u} = (u_x, u_y)$ προστίθεται στην τρέχουσα εκτίμηση: $\mathbf{d}_i \leftarrow \mathbf{d}_i + \mathbf{u}$.

    \item \textbf{Κριτήριο Σύγκλισης:} Ο βρόχος τερματίζει πρόωρα όταν $\max_i(|u_x^i|, |u_y^i|) < 0.01$, υποδηλώνοντας ότι οι διορθώσεις έχουν συγκλίνει. Άνω όριο 100 επαναλήψεων αποτρέπει άπειρο βρόχο σε ιδιαίτερες περιπτώσεις.
\end{enumerate}

\subsubsection*{Δοκιμή σε Τεχνητό Παράδειγμα}

\hspace{0.6cm}Πριν εφαρμοστεί ο αλγόριθμος σε πραγματικά δεδομένα, ελέγχεται η ορθότητά του σε ελεγχόμενες συνθήκες: η $I_2$ δημιουργείται ως καθαρή μετάθεση της $I_1$ κατά $(t_x, t_y)$ pixels μέσω \texttt{cv2.warpAffine}, οπότε η αναμενόμενη ροή είναι γνωστή εκ των προτέρων. Χρησιμοποιείται το crop της αριστερής παλάμης ως περιοχή δοκιμής, με 30 σημεία Shi-Tomasi. Εκτελούνται δύο πειράματα: (α) εναλλαγή $\rho \in \{0.5, 1.0, 3.0, 5.0, 10.0\}$ με σταθερό $\varepsilon=0.005$ και (β) εναλλαγή $\varepsilon \in \{10^{-4}, 10^{-3}, 0.005, 0.05, 0.1, 0.5, 1.0\}$ με σταθερό $\rho=3.0$.

\begin{table}[H]
\centering
\small
\begin{tabular}{@{}cccc@{}}
\toprule
\textbf{Μετάθεση $(t_{x}, t_{y})$} & \textbf{$\rho$} & \textbf{Εκτ. $(u_x,\,u_y)$} & \textbf{Σφάλμα (px)} \\
\midrule
$(1,\,0)$ & 0.5  & $(+1.00,\;+0.07)$ & 0.070 \\
          & 3.0  & $(+0.94,\;-0.01)$ & 0.059 \\
          & 10.0 & $(+0.94,\;-0.01)$ & 0.058 \\
\midrule
$(2,\,0)$ & 0.5  & $(+1.59,\;-0.06)$ & 0.414 \\
          & 3.0  & $(+1.94,\;-0.02)$ & 0.063 \\
          & 10.0 & $(+1.93,\;-0.01)$ & 0.069 \\
\midrule
$(5,\,0)$ & 0.5  & $(+2.02,\;-1.55)$ & 3.359 \\
          & 3.0  & $(+4.14,\;-0.48)$ & 0.989 \\
          & 10.0 & $(+4.97,\;-0.01)$ & 0.032 \\
\midrule
$(5,\,5)$ & 0.5  & $(+0.64,\;-1.66)$ & 7.955 \\
          & 3.0  & $(+3.40,\;+1.13)$ & 4.184 \\
          & 10.0 & $(+4.99,\;+5.00)$ & 0.007 \\
\bottomrule
\end{tabular}
\caption{Πείραμα 1. Επίδραση $\rho$ (σταθερό $\varepsilon=0.005$). Σφάλμα = ευκλείδεια απόσταση από αναμενόμενη ροή.}
\label{tab:lk_rho}
\end{table}

\textbf{Επίδραση $\rho$:}
\begin{itemize}
    \item \textbf{Μικρές μεταθέσεις ($t=1$ px):} Όλες οι τιμές $\rho$ δίνουν παρόμοιο σφάλμα ($\approx 0.06$--$0.07$ px). Ο αλγόριθμος συγκλίνει σωστά ανεξαρτήτως παραθύρου, επιβεβαιώνοντας την ορθότητα της υλοποίησης.

    \item \textbf{Μεγαλύτερες μεταθέσεις ($t_x = 2$--$5$ px):} Το σφάλμα αρχίζει και εξαρτάται έντονα από το $\rho$. Πλέον έχουμε αποδεκτά αποτελέσματα μόνο για μεγάλες τιμές του $\rho$. Αυτό είναι λογικό γιατί το μέγεθος του παραθύρου ολοκλήρωσης πρέπει να είναι ανάλογα μεγάλο της κίνησης που εκτελείται. Ωστόσο παρατηρούμε πως για όλες τις κινήσεις (ακόμα και τις πολύ μεγάλες που γνωρίζουμε ότι ο αλγόριθμος θα έπρεπε να αποτύχει) μεγαλύτερο ρ οδηγεί σε σταδιακά καλύτερα αποτελέσματα έως ότου αυτά συγκλίνουν στα "ιδανικά". Αυτό συμβαίνει ακριβώς επειδή η βασική υπόθεση του LK για σταθερή μετατόπιση εντός του παραθύρου ολοκλήρωσης ισχύει από την στιγμή που όλη η εικόνα έχει την ίδια μετατόπιση. Επομένως κατάλληλα μεγάλο παράθυρο είναι ικανό να ανιχνεύσει την κίνηση και όσο πιο μεγάλο είναι αυτό τόσο πιο πολύ ορθή πληροφορία προσφέρεται κατά την κατασκευή του δομικού τανυστή ο οποίος θα οδηγεί σε πιο καλές εκτιμήσεις.
\end{itemize}


\begin{table}[H]
\centering
\small
\begin{tabular}{@{}cccc@{}}
\toprule
\textbf{Μετάθεση $(t_{x}, t_{y})$} & \textbf{$\varepsilon$} & \textbf{Εκτ. $(u_x,\,u_y)$} & \textbf{Σφάλμα (px)} \\
\midrule
$(1,\,0)$ & $10^{-4}$ & $(+0.946,\;-0.006)$ & 0.055 \\
          & $0.005$   & $(+0.942,\;-0.009)$ & 0.059 \\
          & $0.1$     & $(+0.922,\;-0.046)$ & 0.091 \\
          & $1.0$     & $(+0.747,\;-0.135)$ & 0.286 \\
\midrule
$(2,\,0)$ & $10^{-4}$ & $(+1.946,\;-0.020)$ & 0.058 \\
          & $0.005$   & $(+1.942,\;-0.024)$ & 0.063 \\
          & $0.1$     & $(+1.952,\;-0.043)$ & 0.065 \\
          & $1.0$     & $(+1.740,\;-0.192)$ & 0.323 \\
\midrule
$(5,\,0)$ & $10^{-4}$ & $(+4.141,\;-0.480)$ & 0.984 \\
          & $0.005$   & $(+4.138,\;-0.485)$ & 0.989 \\
          & $0.05$    & $(+4.637,\;-0.825)$ & 0.901 \\
          & $1.0$     & $(+3.508,\;-0.689)$ & 1.644 \\
\midrule
$(5,\,5)$ & $10^{-4}$ & $(+3.331,\;+1.061)$ & 4.278 \\
          & $0.005$   & $(+3.403,\;+1.132)$ & 4.184 \\
          & $0.05$    & $(+3.464,\;+1.615)$ & 3.718 \\
          & $1.0$     & $(+1.857,\;-0.419)$ & 6.265 \\
\bottomrule
\end{tabular}
\caption{Πείραμα 2. Επίδραση $\varepsilon$ (σταθερό $\rho=3.0$).}
\label{tab:lk_eps}
\end{table}


\textbf{Επίδραση $\varepsilon$:}
\begin{itemize}
    \item \textbf{Μικρό $\varepsilon$ ($10^{-4}$--$0.005$):} Για μικρές μεταθέσεις ($t=1,2$ px) τα αποτελέσματα είναι σχεδόν πανομοιότυπα. Eφόσον τα σημεία Shi-Tomasi βρίσκονται σε γωνίες, ο $M$ είναι ήδη καλά καθορισμένος και το $\varepsilon$ παίζει ελάχιστο ρόλο.

    \item \textbf{Μέτριο $\varepsilon$ ($0.05$):} Για μεγαλύτερες μεταθέσεις ($t=5$) όπου ο αλγόριθμος δυσκολεύεται, το $\varepsilon=0.05$ δίνει ελαφρώς καλύτερο αποτέλεσμα από πολύ μικρό $\varepsilon$. Tο regularization σταθεροποιεί αριθμητικά τον $M$ σε περιπτώσεις όπου η εκτίμηση βρίσκεται σε ασταθή περιοχή.

    \item \textbf{Μεγάλο $\varepsilon$ ($\geq 0.5$):} Σαφής υποβάθμιση σε όλες τις μεταθέσεις. Για $\varepsilon=1.0$ και $t_x=1$, η εκτίμηση πέφτει σε $u_x=0.75$ αντί για $1.0$. Tο μεγάλο $\varepsilon$ κυριαρχεί στην διαγώνιο του $M$ και συρρικνώνει τη λύση συστηματικά προς το μηδέν. Η επιλεγμένη τιμή $\varepsilon=0.005$ αποτελεί καλή ισορροπία σταθερότητας και ακρίβειας.
\end{itemize}

\subsection*{Σύγκριση LK και TV-L1}

\hspace{0.6cm}Για τον υπολογισμό και σύγκριση της ροής σε τυχαία ζεύγη πλαισίων υλοποιήθηκε η συνάρτηση \texttt{compare\_lk\_vs\_tvl1()}. Υπάρχει ένας θεμελιώδης σχεδιαστικός περιορισμός: για ένα τυχαίο ζεύγος πλαισίων $(I_{t-1}, I_{t})$ δεν είναι γνωστή η θέση των bounding boxes καθώς αυτή προκύπτει μόνο μετά την εκτέλεση του tracking. Για να αποφευχθεί αυτή η εξάρτηση, εντοπίζονται $N=90$ σημεία Shi-Tomasi σε ολόκληρη την εικόνα (περίπου 30 ανά υποτιθέμενο box), χωρίς περιορισμό σε συγκεκριμένη περιοχή. Η ροή υπολογίζεται στα σημεία αυτά και οι δύο μέθοδοι συγκρίνονται ισότιμα.

\hspace{0.6cm}Η σύγκριση πραγματοποιήθηκε σε 6 ζεύγη πλαισίων (frames 0-1, 10-11, 35-36, 40-41, 50-51, 66-67), που καλύπτουν σκηνές από σχεδόν ακίνητες έως έντονης κίνησης χεριών. Σε κάθε ζεύγος εντοπίστηκαν $N=90$ σημεία Shi-Tomasi στο $I_n$, υπολογίστηκε ροή με LK και πυκνή ροή με TV-L1, ενώ το TV-L1 πεδίο δειγματοληπτήθηκε στις θέσεις $\mathbf{x}_i + \mathbf{d}_i^{\text{LK}}$ (από το σημείο δηλαδή που ο LK προβλέπει ότι ήρθαν τα σημεία ενδιαφέροντος της εικόνας $I_{n}$) για αντιστοιχία. Ως μέτρο σύγκρισης χρησιμοποιείται η μέση ευκλείδεια απόσταση $\overline{D} = \frac{1}{N}\sum_i \|\mathbf{d}_i^{\text{LK}} - \mathbf{d}_i^{\text{TVL1}}\|_2$, όπου i το i-οστό feature.

\begin{figure}[H]
    \centering
    \begin{subfigure}[b]{0.32\textwidth}
        \centering
        \includegraphics[width=\linewidth]{lk-tvl1_comparison/comp_000.png}
        \caption{Ζεύγος 0 ($\overline{D}=0.119$)}
    \end{subfigure}
    \hfill
    \begin{subfigure}[b]{0.32\textwidth}
        \centering
        \includegraphics[width=\linewidth]{lk-tvl1_comparison/comp_001.png}
        \caption{Ζεύγος 1 ($\overline{D}=0.359$)}
    \end{subfigure}
    \hfill
    \begin{subfigure}[b]{0.32\textwidth}
        \centering
        \includegraphics[width=\linewidth]{lk-tvl1_comparison/comp_002.png}
        \caption{Ζεύγος 2 ($\overline{D}=0.207$)}
    \end{subfigure}

    \vspace{0.3cm}

    \begin{subfigure}[b]{0.32\textwidth}
        \centering
        \includegraphics[width=\linewidth]{lk-tvl1_comparison/comp_003.png}
        \caption{Ζεύγος 3 ($\overline{D}=0.565$)}
    \end{subfigure}
    \hfill
    \begin{subfigure}[b]{0.32\textwidth}
        \centering
        \includegraphics[width=\linewidth]{lk-tvl1_comparison/comp_004.png}
        \caption{Ζεύγος 4 ($\overline{D}=0.160$)}
    \end{subfigure}
    \hfill
    \begin{subfigure}[b]{0.32\textwidth}
        \centering
        \includegraphics[width=\linewidth]{lk-tvl1_comparison/comp_005.png}
        \caption{Ζεύγος 5 ($\overline{D}=0.696$)}
    \end{subfigure}

    \caption{Σύγκριση LK (κόκκινο) και TV-L1 (κυανό) για 6 ζεύγη πλαισίων. Τα βέλη αναπαριστούν forward ροή, τοποθετoύνται στη εικόνα $I_{n-1}$ και αναπαριστούν την πρόβλεψη της κίνησης.}
    \label{fig:lk_tvl1_comparison}
\end{figure}

\subsection*{Παρατηρήσεις Αποτελεσμάτων:}

\begin{itemize}
    \item \textbf{Συμφωνία σε Μικρές Μετατοπίσεις:} Στα ζεύγη 0 και 4 ($\overline{D} < 0.17$), όπου η σκηνή είναι σχεδόν ακίνητη, οι δύο μέθοδοι συμφωνούν πλήρως. Η γραμμική υπόθεση Taylor που εφαρμόζει ο LK ισχύει για μικρές μετατοπίσεις, οπότε καταλήγει στην ίδια λύση με το TV-L1.

    \item \textbf{Αποκλίσεις σε Έντονη Κίνηση:} Στα ζεύγη 3 και 5 ($\overline{D} = 0.565$ και $0.696$), όπου τα χέρια εκτελούν ταχείες χειρονομίες, η απόκλιση αυξάνεται σημαντικά. Ο LK λόγω της γραμμικής προσέγγισης αποτυγχάνει όταν η πραγματική μετατόπιση υπερβαίνει το μέγεθος του παραθύρου παρατήρησης και η υπόθεση $I_{n-1}(\mathbf{x}+\mathbf{d}) \approx I_n(\mathbf{x}) + \nabla I \cdot \mathbf{d}$ παύει να ισχύει. Το TV-L1 αντιμετωπίζει αυτό το πρόβλημα μέσω κλιμακωτής ανάλυσης εσωτερικά, επιτυγχάνοντας αξιόπιστα αποτελέσματα και σε μεγάλες μετατοπίσεις.

    \item \textbf{Κατεύθυνση vs. Μέγεθος:} Σε αρκετά ζεύγη (π.χ. ζεύγος 2), οι δύο μέθοδοι συμφωνούν στη γενική κατεύθυνση της κίνησης (π.χ. χέρι προς το σώμα) αλλά διαφέρουν στο μέγεθος των διανυσμάτων. Ο TV-L1 παράγει χωρικά ομογενή πεδία λόγω του regularization όρου που τιμωρεί τις απότομες μεταβολές στο πεδίο της ροής και επιβάλλει ομαλότητα μεταξύ γειτονικών pixels. Ο LK αντίθετα εκτιμά κάθε σημείο ανεξάρτητα, χωρίς χωρική κανονικοποίηση, με αποτέλεσμα πιο ανομοιόμορφα διανύσματα μεταξύ γειτονικών features.

    \item \textbf{Αραιή vs. Πυκνή Ροή:} Ο LK υπολογίζει ροή μόνο στα $N$ σημεία Shi-Tomasi (90 ανά ζεύγος), ενώ το TV-L1 παράγει πυκνό πεδίο ροής σε κάθε pixel. Αυτό καθιστά τον LK κατάλληλο για εφαρμογές tracking σε συγκεκριμένα σημεία (π.χ. κεφάλι, χέρια), ενώ το TV-L1 είναι απαραίτητο για εφαρμογές που χρειάζονται πλήρη πεδία ροής.

    \item \textbf{Υπολογιστικό Κόστος:} Ο LK επιλύει ένα 2×2 σύστημα ανά σημείο σε λίγες επαναλήψεις και είναι κατάλληλος για real-time. Το TV-L1 απαιτεί 1-2 τάξεις μεγέθους περισσότερο χρόνο επεξεργασίας για ένα ζεύγος πλαισίων.
\end{itemize}

\subsubsection*{Επίδραση Παραμέτρων $\rho$ και $\varepsilon$ στη Σύγκριση LK--TV-L1}

\hspace{0.6cm}Για να αξιολογηθεί πώς επηρεάζουν οι παράμετροι του LK τη συμφωνία του με το TV-L1, επαναλήφθηκε η σύγκριση για όλους τους συνδυασμούς $\rho \in \{0.5,\,1.0,\,2.0,\,5.0,\,10.0\}$ και $\varepsilon \in \{0,\,0.01,\,0.05,\,0.1,\,1.0\}$. Ο Πίνακας~\ref{tab:lk_tvl1_params} παρουσιάζει τη μέση απόσταση $\overline{D}$ (μέσος όρος επί 6 ζευγών πλαισίων).

\begin{table}[H]
\centering
\small
\begin{tabular}{@{}c ccccc@{}}
\toprule
 & \multicolumn{5}{c}{\textbf{$\varepsilon$}} \\
\textbf{$\rho$} & $0$ & $0.01$ & $0.05$ & $0.1$ & $1.0$ \\
\midrule
$0.5$  & 3.32 & 1.39 & 1.41 & 1.39 & 1.46 \\
$1.0$  & 2.83 & 1.27 & 1.27 & 1.22 & 1.21 \\
$2.0$  & 0.66 & 0.66 & 0.64 & 0.61 & 0.91 \\
$5.0$  & 0.42 & 0.42 & 0.38 & \textbf{0.34} & 0.67 \\
$10.0$ & 0.47 & 0.46 & 0.43 & 0.41 & 0.74 \\
\bottomrule
\end{tabular}
\caption{Μέση απόσταση $\overline{D}$ LK--TV-L1 ανά παράμετρο (μέσος όρος 6 ζευγών πλαισίων). Bold: η καλύτερη τιμή.}
\label{tab:lk_tvl1_params}
\end{table}

\textbf{Παρατηρήσεις:}
\begin{itemize}
    \item \textbf{Επίδραση $\rho$:} Η αύξηση του $\rho$ βελτιώνει σταθερά (για όλα τα epsilon) τη συμφωνία με το TV-L1 μέχρι $\rho=5.0$ ($\overline{D}=0.34$). Το $\rho=10.0$ δίνει ελαφρώς χειρότερο αποτέλεσμα ($\overline{D}=0.41$) παρότι στο toy example ήταν το καλύτερο. Η αιτία όπως αναφέρθηκε είναι ότι εδώ τα 90 σημεία εντοπίζονται σε \textit{ολόκληρη} την εικόνα, καλύπτοντας χέρια, πρόσωπο και φόντο που κινούνται διαφορετικά. Με $\rho=10.0$ το παράθυρο ($61\times61$ pixels, $\sim$25\% της εικόνας) αναμιγνύει κινήσεις από διαφορετικές περιοχές, ενώ το TV-L1 εκτιμά τη ροή ανά pixel συνεπώς η εκτίμηση αποκλίνει. Αξίζει να σημειωθεί ωστόσο ότι μεγαλύτερα ρ έχουν καλύτερα αποτελέσματα σε μεγαλύτερες κινήσεις και χειρότερα για μικρότερες κινήσεις. Η αξιολόγηση που γίνεται με τον μέσο όρο δεν απόλυτα αντιπροσωπευτική και η εποπτεία των συνολικών αποτελεσμάτων που βρίσκονται στον κώδικα είναι επιβεβλήμένη για την εξαγωγή ασφαλών συμπερασμάτων, με δεδομένο ότι τα ζεύγη επιλέχθηκαν εντελώς τυχαία και μπορεί να περιέχουν περισσότερες "μικρές" από ότι "μεγάλες" κινήσεις.

    \item \textbf{Επίδραση $\varepsilon$ :} Για μικρό $\rho$ ($0.5$, $1.0$) με $\varepsilon=0$, οι αποστάσεις εκτοξεύονται ($3.32$ και $2.83$) λόγω αριθμητικής αστάθειας. Ακόμα και μικρό $\varepsilon=0.01$ μειώνει δραστικά την απόσταση (π.χ.\ από $8.70$ σε $2.90$ για ζεύγος~1 με $\rho=0.5$). Για μεγάλο $\rho$, το $\varepsilon$ έχει μικρότερη επίδραση γιατί ο $M$ είναι ήδη καλά καθορισμένος. Μεγάλο $\varepsilon$ ($=1.0$) χειροτερεύει το αποτέλεσμα σε όλες τις περιπτώσεις αφού έτσι τα διανύσματα LK τείνουν στο μηδέν.

\end{itemize}
% ============================================================
\subsubsection*{1.1.2 Υπολογισμός της Μετατόπισης των Παραθύρων από τα Διανύσματα Οπτικής
Ροής}
% ============================================================

\hspace{0.6cm}Στο ερώτημα αυτό εφαρμόζεται tracking τριών bounding boxes (πρόσωπο, αριστερό χέρι, δεξί χέρι) σε ακολουθία 70 πλαισίων. Σε κάθε βήμα, τα όρια του box στην $Ι_{n-1}$ χρησιμοποιούνται για τον ορισμό του crop στην $I_{n}$ εικόνα από το οποίο θα προκύψουν 30 σημείων Shi-Tomasi. Στη συνέχεια εκτιμάται η οπτική ροή με LK (ή πυκνή ροή με TV-L1), και η μετατόπιση κάθε box υπολογίζεται από τη συνάρτηση \texttt{displ()} με δύο κριτήρια: (α) \textbf{κατώφλι ενέργειας}, επιλέγοντας τα $d_{i}$ των οποίων η ενέργεια ξεπερνά ένα κατώφλι και υπολογίζοντας τον μέσο όρο τους, και (β) \textbf{διάμεσος}, υπολογίζοντας την διάμεσο των $d_{i}$ ανεξαρτήτως ενέργειας. Τα σφάλματα εκτίμησης συσσωρεύονται διαδοχικά από πλαίσιο σε πλαίσιο.

\subsubsection*{Επίδραση της Παραμέτρου $\rho$}

\hspace{0.6cm}Η παράμετρος $\rho$ καθορίζει το μέγεθος του παραθύρου ολοκλήρωσης και κατά συνέπεια την ευστάθεια των εκτιμήσεων ροής. Η ακόλουθη εικόνα (Σχήμα~\ref{fig:tracking_rho}) συγκρίνει τα αποτελέσματα tracking για $\rho \in \{0.5, 3.0, 10.0\}$ (με $\varepsilon=0$) στα πλαίσια 35 και 60.



\begin{itemize}
    \item \textbf{Μικρό $\rho=0.5$:} Το παράθυρο ολοκλήρωσης ($4\times4$ pixels) είναι υπερβολικά μικρό. Ο τανυστής $M$ δεν λαμβάνει επαρκή πληροφορία, με αποτέλεσμα να κάνει θορυβώδεις εκτιμήσεις ροής (τα βέλη στο πλαίσιο 35 είναι ακανόνιστα) και να χάνει τα boxes ήδη από το πλαίσιο 35.
    
    \item \textbf{Μέτριο $\rho=3.0$:} Το παράθυρο ($19\times19$ pixels) εξισορροπεί τοπικότητα και σταθερότητα. Τα βέλη είναι ομοιόμορφα και σε σωστές κατευθύνσεις, και τα boxes παραμένουν ευθυγραμμισμένα.

\begin{figure}[H]
    \centering
    \begin{subfigure}[b]{0.32\textwidth}
        \centering
        \includegraphics[width=\linewidth]{track_lk/rho_0.5_eps_0/frame_035.png}
        \caption{$\rho=0.5$, πλαίσιο 35}
    \end{subfigure}
    \hfill
    \begin{subfigure}[b]{0.32\textwidth}
        \centering
        \includegraphics[width=\linewidth]{track_lk/rho_3.0_eps_0.05/frame_035.png}
        \caption{$\rho=3.0$, πλαίσιο 35}
    \end{subfigure}
    \hfill
    \begin{subfigure}[b]{0.32\textwidth}
        \centering
        \includegraphics[width=\linewidth]{track_lk/rho_10_eps_0/frame_035.png}
        \caption{$\rho=10.0$, πλαίσιο 35}
    \end{subfigure}

    \vspace{0.2cm}

    \begin{subfigure}[b]{0.32\textwidth}
        \centering
        \includegraphics[width=\linewidth]{track_lk/rho_0.5_eps_0/frame_060.png}
        \caption{$\rho=0.5$, πλαίσιο 60}
    \end{subfigure}
    \hfill
    \begin{subfigure}[b]{0.32\textwidth}
        \centering
        \includegraphics[width=\linewidth]{track_lk/rho_3.0_eps_0.05/frame_060.png}
        \caption{$\rho=3.0$, πλαίσιο 60}
    \end{subfigure}
    \hfill
    \begin{subfigure}[b]{0.32\textwidth}
        \centering
        \includegraphics[width=\linewidth]{track_lk/rho_10_eps_0/frame_060.png}
        \caption{$\rho=10.0$, πλαίσιο 60}
    \end{subfigure}

    \caption{Επίδραση $\rho$ στο tracking: $\rho=0.5$ (αριστερά), $\rho=3.0$ (κέντρο), $\rho=10.0$ (δεξιά), πλαίσια 35 και 60.}
    \label{fig:tracking_rho}
\end{figure}

    \item \textbf{Μεγάλο $\rho=10.0$:} Το παράθυρο ($61\times61$ pixels) ενσωματώνει ευρύτερη γειτονιά. Αυτό θα περιμέναμε ναι μεν να σταθεροποιεί τα βέλη σε σκηνές ομοιόμορφης κίνησης, αλλά να δημιουργεί πρόβλημα σε ταχείες χειρονομίες καθώς το μεγάλο παράθυρο αναμιγνύει πληροφορία από διαφορετικές περιοχές (χέρι/φόντο), οδηγώντας σε λανθασμένη εκτίμηση μετατόπισης. Κάτι τέτοιο ωστόσο δεν παρατηρείται στα πειραματικά αποτελέσματα. Μια ερμηνεία που μπορεί αν δωθεί είναι πως ακριβώς επειδή εργαζόμαστε σε αποκομμένα boxes (που οριοθετούν τα μέλη που μας αφορούν), όσο μεγάλο και αν είναι το παράθυρο ολοκλήρωσης δεν πρόκειται να συμπεριληφθεί άλλο μέλος (π.χ να μπει στο παράθυρο του χεριού το πρόσωπο) ούτε και το φόντο. Το παράθυρο θα λαμβάνει ακόμα περισσότερη πληροφορία από το μέλος που το αφορά (και μόνο από αυτό) αγνοώντας μάλιστα τις πολύ λεπτομερείς κινήσεις που ενδεχομένως δεν συνεισφέρουν στην συνολική (μεταφορική) κίνηση αλλά περιγράφουν εσωτερικές αναδιατάξεις (π.χ κούνημα δακτύλων με σταθερή παλάμη ή μετακίνηση βλεφάρων και στόματος με σταθερό το υπόλοιπο κεφάλι). Αυτό οδηγεί σε καλύτερο συνολικό tracking το οποίο φαίνεται ξεκάθαρα από τις παραπάνω εικόνες.
\end{itemize}

\subsubsection*{Επίδραση της Παραμέτρου $\varepsilon$ στο Tracking}

\hspace{0.6cm}Η επίδραση του $\varepsilon$ στο tracking ακολουθεί τα ίδια μοτίβα που παρατηρήθηκαν στο τεχνητό παράδειγμα:

\begin{itemize}
    \item \textbf{$\varepsilon=0$ με μικρό $\rho$:} Ο τανυστής $M$ μπορεί να είναι σχεδόν ιδιόμορφος (σε ομοιογενείς περιοχές), οδηγώντας σε αριθμητικά ασταθείς εκτιμήσεις.

    \item \textbf{Μεγάλο $\varepsilon=1.0$:} Συρρικνώνει τα διανύσματα μετατόπισης προς το μηδέν. Συνεπώς εμφανίζεται μια αδράνεια κατά την μετακίνηση των boxes τα οποία ουσιαστικά παραμένουν στη θέση τους, χωρίς να παρακολουθούν την πραγματική κίνηση

    \item \textbf{Βέλτιστο $\varepsilon \in [0.01, 0.1]$:} Παρέχει αριθμητική σταθεροποίηση χωρίς υπερβολική συρρίκνωση, επιτρέποντας στον αλγόριθμο να αναγνωρίσει αξιόπιστα τις κατευθύνσεις κίνησης.
\end{itemize}

\subsubsection*{Επίδραση Κριτηρίου Εκτίμησης Μετατόπισης: Κατώφλι vs.\ Διάμεσος}

\hspace{0.6cm} Η σύγκριση πραγματοποιήθηκε στο πλαίσιο 60:

\begin{figure}[H]
    \centering
    \begin{subfigure}[b]{0.32\textwidth}
        \centering
        \includegraphics[width=\linewidth]{track_lk/threshold_0.001/frame_060.png}
        \caption{Κατώφλι = 0.001}
    \end{subfigure}
    \hfill
    \begin{subfigure}[b]{0.32\textwidth}
        \centering
        \includegraphics[width=\linewidth]{track_lk/threshold_0.5/frame_060.png}
        \caption{Κατώφλι = 0.5}
    \end{subfigure}
    \hfill
    \begin{subfigure}[b]{0.32\textwidth}
        \centering
        \includegraphics[width=\linewidth]{track_lk/median/frame_060.png}
        \caption{Διάμεσος}
    \end{subfigure}
    \caption{Σύγκριση κριτηρίων εκτίμησης μετατόπισης στο πλαίσιο 60. Αριστερά: χαμηλό κατώφλι (0.001) που επιλέγει πολλά features· κέντρο: υψηλό κατώφλι (0.5) που επιλέγει μόνο τα ισχυρά· δεξιά: χρήση διαμέσου.}
    \label{fig:tracking_displ}
\end{figure}

\begin{itemize}
    \item \textbf{Χαμηλό κατώφλι (thr $= 0.001$):} Λαμβάνει υπόψην για τον υπολογοσμό της συνολικής μετατόπισης και εσωτερικές αναδιατάξεις (μικρές κινήσεις) αλλά και θορυβώδεις τιμές ροής με αποτέλεσμα ο μέσος όρος να μολύνεται και να έχουμε ελαφρά αστάθεια στο tracking.

    \item \textbf{Υψηλό κατώφλι (thr $= 0.5$--$1.0$):} Επιλέγει μόνο τις ισχυρές μετατοπίσεις. Αυτό βελτιώνει την ακρίβεια ανά βήμα αλλά μπορεί για μεγαλύτερες τιμές (π.χ >3) να αφήσει ελάχιστες ροές να αποφασίσουν για την μετατόπιση, ροές που μάλιστα με τόσο μεγάλη ενέργεια πιθανόν είναι και λανθασμένες (αφού γνωρίοζουμε οτι οι αλγόριθμοι δεν αποδίδουν σε μαγάλες μετατοπίσεις). Για ακραία μεγάλες τιμές προφανώς κανένα διάνυσμα δεν περνάει και το box παραμένει ακίνητο.

    \item \textbf{Διάμεσος:} Είναι στατιστικά ανθεκτικός καθώς μεμονωμένα features με λανθασμένη ροή δεν επηρεάζουν σημαντικά τον διάμεσο. Από την άλλη εξασφαλίζει ότι θα υπάρχει και ένα καλό δείγμα από το οποίο θα προκύψει η συνολικη μετατόπιση (τα μισά).
\end{itemize}

Το μειονέκτημα των παραπάνω μεθόδων είναι ότι αγνοούν την ποιότητα των features και δίνουν ίσο βάρος σε ισχυρές και αδύναμες γωνίες. Επίσης ενδεχομένως και ένα άνω όριο ενέργειας (upper threshold) να αποτελούσε βελτίωση της απλής μεθόδου με την λογική της απόρριψης ακραίων μετατοπίσεων που έτσι και αλλιώς δεν μπορούν να ανιχνευθούν. Προφανώς η παρατήρηση σε δράσεις που εκτυλίσσονται σε πολλά περισσότερα frames θα αναδείκνυαν και πολύ σημαντικότερες διαφορές ανάμεσα στις μεταβολές των παραμέτρων και μεθόδων.

\subsubsection*{Σύγκριση Tracking LK και TV-L1}

\hspace{0.6cm}Το TV-L1 χρησιμοποιεί πυκνή οπτική ροή (ανά pixel) για τον υπολογισμό της μετατόπισης των boxes. Αντί να εντοπίζει 30 σημεία Shi-Tomasi, δειγματοληπτεί τη ροή σε πλέγμα πυκνών σημείων εντός κάθε crop, εξαλείφοντας την εξάρτηση από την ανίχνευση γωνιών. Η σύγκριση στα πλαίσια 35 και 69 παρουσιάζεται παρακάτω:

\begin{figure}[H]
    \centering
    \begin{subfigure}[b]{0.48\textwidth}
        \centering
        \includegraphics[width=\linewidth]{track_lk/rho_3.0_eps_0.05/frame_035.png}
        \caption{LK ($\rho=3.0,\,\varepsilon=0.05$), πλαίσιο 35}
    \end{subfigure}
    \hfill
    \begin{subfigure}[b]{0.48\textwidth}
        \centering
        \includegraphics[width=\linewidth]{track_tvl1/frame_035.png}
        \caption{TV-L1, πλαίσιο 35}
    \end{subfigure}

    \vspace{0.2cm}

    \begin{subfigure}[b]{0.48\textwidth}
        \centering
        \includegraphics[width=\linewidth]{track_lk/rho_3.0_eps_0.05/frame_069.png}
        \caption{LK ($\rho=3.0,\,\varepsilon=0.05$), πλαίσιο 69}
    \end{subfigure}
    \hfill
    \begin{subfigure}[b]{0.48\textwidth}
        \centering
        \includegraphics[width=\linewidth]{track_tvl1/frame_069.png}
        \caption{TV-L1, πλαίσιο 69}
    \end{subfigure}

    \caption{Σύγκριση tracking LK (αριστερά) και TV-L1 (δεξιά) στα πλαίσια 35 (πάνω) και 69 (κάτω).}
    \label{fig:tracking_lk_tvl1}
\end{figure}

Και οι δύο μέθοδοι παρακολουθούν αξιόπιστα πρόσωπο και χέρια, με παρόμοια θέση boxes. Η ταχύτητα LK είναι σημαντικά μεγαλύτερη γιατί επιλύει 2×2 συστήματα μόνο για 30 features ανά box, ενώ το TV-L1 υπολογίζει πυκνό πεδίο για ολόκληρο το πλαίσιο.

Ο TV-L1 παρουσιάζει πιο ομαλές ροές στα σημεία ενδιαφέροντος γεγονός που οφείλεται στον regularization όρο που τιμωρεί τις ασυνέχειες, συνεπώς πειρμένουμε να προβλέπει μια πιο αναμενόμενη κίνηση. Ωστόσο και οι δύο μέθοδοι στο σύντομο αυτό χρονικό διάστημα παρουσιάζουν παραπλήσιες ικανοποιητικές αποδόσεις και αποτυγχάνουν ταυτόγχρονα στα τελικά frames όπου η κίνηση σε διαδοχικά frames είναι πολύ μεγάλη, σπάζοντας τις υποθέσεις των αλγορίθμων.


% ============================================================
\subsubsection*{1.1.3 Πολυκλιμακωτή LK (Multiscale Lucas-Kanade)}
% ============================================================

\subsubsection*{Θεωρητική Βάση}

\hspace{0.6cm}Ο βασικός περιορισμός του LK είναι η γραμμική προσέγγιση της εξίσωσης φωτεινότητας:
\[
    I_{n-1}(\mathbf{x} + \mathbf{d}) \approx I_{n-1}(\mathbf{x}) + \nabla I_{n-1}(\mathbf{x})^\top \mathbf{d}
\]
Η προσέγγιση αυτή ισχύει μόνο για μικρές μετατοπίσεις $\|\mathbf{d}\|$, τυπικά $\leq 1$--$2$ pixels ανάλογα με το $\rho$. Για μεγάλες κινήσεις χεριών (π.χ.\ $5$--$20$ pixels/πλαίσιο) ο αλγόριθμος επιστρέφει λανθασμένη ροή.

Η λύση είναι η πολυκλιμακωτή εκδοχή του αλγορίθμου LK. Για την δημιουργία μιας πυραμίδας από εικόνες κατασκευάζονται n\_scales επίπεδα. Κάθε επίπεδο προκύπτει από το προηγούμενο με εφαρμογή gaussian φίλτρου ακολουθούμενη από υποδειγματοληψία του αποτελέσματος κατά παράγοντα 2. Το φιλτράρισμα με gaussian αποτρέπει το φαινόμενο aliasing κατά την υποδειγματοληψία. Η εκτίμηση της ροής d διαδίδεται από την κορυφή της πυραμίδας προς τα κάτω με διπλασιασμό ώστε να ταιριάξει στις νέες διαστάσεις της εικόνας.

Κάθε επίπεδο επιλύει ένα πρόβλημα μικρής μετατόπισης ακόμα και αν η συνολική κίνηση είναι μεγάλη.

\subsubsection*{Υλοποίηση \texttt{multi\_scale\_lk()}}

\hspace{0.6cm}Σημαντικά σημεία της συνάρτησης \texttt{multi\_scale\_lk(I1, I2, features, $\rho$, $\varepsilon$, n\_scales)} είναι:

\begin{enumerate}

    \item \textbf{Εκκίνηση από ανώτερο επίπεδο:} Αρχικοποίηση $\mathbf{d}^{(L-1)} = \mathbf{0}$. Τα σημεία scaled: $\mathbf{x}^{(s)} = \mathbf{x} / 2^s$.

    \item \textbf{Επανάληψη LK ανά επίπεδο:} Σε κάθε $s$ (από $L{-}1$ έως $0$) καλείται \texttt{lk()} με αρχικοποίηση από το προηγούμενο επίπεδο μέσω των προαιρετικών ορισμάτων \texttt{d\_x0, d\_y0}.

    \item \textbf{Έξοδος:} Το τελικό $\mathbf{d}^{(0)}$ αντιστοιχεί σε πλήρη ανάλυση.
\end{enumerate}

\subsubsection*{Σύγκριση Tracking - Πολυκλιμακωτός LK vs LK}

\hspace{0.6cm} Στην παρακάτω εικόνα (figure 5) αντιπαραβάλλονται τα αποτελέσματα των δύο αλγορίθμων στα frames 47 και 69. Το αποτέλεσμα του tracking μέσω του πολυκλιμακωτού LK είναι ανώτερο. Αυτό γίνεται αντιληπτό ιδιαίτερα από το frame 42 και έπειτα, όπου η κίνηση του χεριού γίνεται πιο απότομη μεταξύ διαδοχικών frames. Το bounding box του αριστερού χεριού, ωστόσο συνεχίζει να περικλείει το χέρι ενώ στον κλασσικό LK έχει ολισθήσει σημαντικά. Η καλύτερη απόκριση του αλγορίθμου πιθανώς οφείλεται: 

\begin{figure}[H]
    \centering
    \begin{subfigure}[b]{0.48\textwidth}
        \centering
        \includegraphics[width=\linewidth]{track_multiscale_lk/frame_047.png}
        \caption{frame 47 - ml\_lk}
    \end{subfigure}
    \hfill
    \begin{subfigure}[b]{0.48\textwidth}
        \centering
        \includegraphics[width=\linewidth]{track_lk/rho_3.0_eps_0.05/frame_047.png}
        \caption{frame 47 - lk}
    \end{subfigure}

    \vspace{0.2cm}

    \begin{subfigure}[b]{0.48\textwidth}
        \centering
        \includegraphics[width=\linewidth]{track_multiscale_lk/frame_069.png}
        \caption{frame 69 - ml\_lk}
    \end{subfigure}
    \hfill
    \begin{subfigure}[b]{0.48\textwidth}
        \centering
        \includegraphics[width=\linewidth]{track_lk/rho_3.0_eps_0.05/frame_069.png}
        \caption{frame 69 - lk}
    \end{subfigure}

    \caption{Multiscale LK tracking ($\rho=3.0$, $\varepsilon=0.05$, $L=4$) vs LK tracking ($\rho=3.0$, $\varepsilon=0.05$).}
    \label{fig:multiscale_tracking}
\end{figure}

\begin{itemize}
    \item \textbf{Χειρισμός μεγάλων μετατοπίσεων:} Η πολυκλιμακωτή προσέγγιση επιτρέπει εκτίμηση μετατοπίσεων έως $2^{L-1}$ φορές μεγαλύτερων με τον ίδιο $\rho$. Το $L=4$ καλύπτει μετατοπίσεις έως $2^3=8$ φορές του αρχικού $\rho$, χωρίς να χρειάζεται αύξηση του $\rho$ (που θα εισήγαγε ανεπιθύμητη ανάμιξη γειτονικών κινήσεων).

    \item \textbf{Καλύτερη αρχικοποίηση:} Αντί για $\mathbf{d}^{(0)} = \mathbf{0}$ (απλός LK), ο multiscale LK ξεκινά κάθε ενδιάμεσο επίπεδο της πυραμίδας από μια ήδη καλή αρχικοποίηση, μειώνοντας τις απαραίτητες επαναλήψεις για σύγκλιση.

    \item \textbf{Σταθερότητα σε ακολουθίες:} Διαδοχικά μικρότερα σφάλματα προφανώς οδηγούν σε μικρότερη ολίσθηση του bounding box από την περιοχή που μας ενδιαφέρει (π.χ αριστερό χέρι). Έτσι η ακολουθία tracking παραμένει ακριβής για περισσότερα πλαίσια.

\end{itemize}

\end{document}
